\chapter{$2$-Power Bernoulli Units, Universal Graded Deformations, and Closure of the Tate--Shafarevich Window Bound}
\label{ch:20}

We now reach the culminating chapter of Part~VI—and the final load-bearing theorem of the entire proof of the Generalized Riemann Hypothesis. In \cref{ch:19} we resolved Strata~0, 1, and~2 of the Bloch--Kato Tate--Shafarevich group $\Sha(E, \ZZ_\ell(e+1))$: the window-avoidance strata ($\ell = 2$, Count~IV, Count~$\mathrm{V}_{\mathrm{sm}}$), the entire odd minus part via Kolyvagin's Gauss-sum Euler system, and the trivial-character outer leaf (Count~Z). What remains is **Stratum~3**: the \emph{odd plus part} ($\chi(c) = (-1)^{e+1}$), which lies in the classic Vandiver-conjecture blind spot where Stickelberger elements vanish identically!

The conceptual breakthrough powering Stratum~3—originating in Ribet's converse to Herbrand's theorem and developed here into a uniform, unconditional quantitative machine across all four plus-part regimes (**Counts~I, II, III, V**)—is to realize every nontrivial plus-part Selmer class as a first-order tangent vector to a \emph{graded reducible deformation ring} of a diagonal residual representation $\bar{r}_0$ into $\mathcal{G}_3 = \GL_3 \rtimes \{\pm 1\}$ (for $\UU(3)$ in Counts~I and~V) or $\GSp_4$ (in Counts~II and~III). Because the diagonal deformation ring admits an exact $X$-graded presentation (\textbf{Lemma~P}), its cotangent length (bounding $\len\Sha^+$ via \textbf{Lemma~M'}) is controlled up to an \emph{exact factor of $2$} (\textbf{Lemma~M}) by the valuation of an Eisenstein congruence element evaluated at the finitely many irreducible cuspidal automorphic points!

This chapter establishes all four pillars of Stratum~3 and executes the final single-pass assembly:
\begin{enumerate}[label=(\roman*)]
\item \emph{$2$-power Bernoulli units and unramified cocycle vanishing} (\cref{thm:20:bernoulli-units}): kills all unwanted off-diagonal root spaces (weight $-2$ in Count~V and all six odd roots $\pm w_{\alpha, \beta, \theta}$ in Counts~II--III) via an elementary $2$-power twist theorem.
\item \emph{Universal graded deformation-ring length bounds (Lemmas~M', M, P with exact factor $2$), unconditional finite-level adjoint rigidity, and one-shot four-line Kato vanishing over $K/\QQ$} (\cref{thm:20:graded-length}): presents the entire commutative-algebraic heart cleanly in optimal logical order ($\mathrm{M'} \to \mathrm{M} \to \mathrm{P}$), proves finite-level adjoint rigidity without big-image hypotheses, and collapses the two-weight Euler-system vanishing to four lines via Kato~\cite[Theorem~14.2(1)]{Kato2004} and Jacquet--Shalika~\cite[Theorem~1]{JacquetShalika1976}.
\item \emph{Unified four-step deformation--congruence machine for Counts~I, II, III, and~V and single-pass $\sqrt{b} = \exp(\frac{1}{2}x^\varepsilon)$ closure} (\cref{thm:20:vwin-complete}): closes all four regimes uniformly (honestly retaining rank-$4$ unitary Hida families plus Ferrero--Washington deep-layer propagation for Counts~II--III), and proves $\mathcal{V}(E, \lfloor x^{401}\rfloor) \le \exp(\frac{1}{2}x^\varepsilon) = \sqrt{b}$ on $[W, W^2]$ in one pass.
\end{enumerate}

\section{Elementary $2$-Power Bernoulli Units and Unramified Cocycle Vanishing}
\label{sec:20:bernoulli}

In Counts~II, III, and~V, the residual diagonal representation $\bar{r}_0$ carries extra root spaces whose Galois cohomology must vanish in order for the graded deformation ring to satisfy the smooth-chart and two-signed-point isolation hypotheses (M1)--(M5) of \cref{thm:20:graded-length}. To kill these auxiliary root spaces unconditionally without assuming Vandiver's conjecture, we twist by an auxiliary primitive Dirichlet character $\xi$ of large $2$-power conductor $2^n$: Washington's theorem on the boundedness of $\ell$-parts of class numbers in $\ZZ_2$-extensions~\cite{Washington1997} implies that the generalized Bernoulli numbers $B_{1, \xi\chi}$ become $\ell$-adic \emph{units} for all sufficiently large $n$!

\litref{Washington (1997, Introduction to Cyclotomic Fields, Theorem~10.4 \& \S7)~\cite{Washington1997}; Ferrero--Washington (1979)~\cite{FerreroWashington1979}}

\begin{theorem}[$2$-power Bernoulli units and unramified cocycle vanishing criterion]
\label{thm:20:bernoulli-units}
Let $\ell \ge 3$ be an odd prime, let $N_0 \ge 1$ be an odd integer, and let $\chi_0$ be a Dirichlet character modulo $N_0$ of prime-to-$2$ conductor (or fixed auxiliary character of conductor dividing $N_h$).
\begin{enumerate}[label=\textup{(\roman*)}]
\item \emph{(Elementary $2$-power Bernoulli unit theorem.)} There exists an explicit integer $n_0 = n_0(\ell, N_0)$ (polynomially bounded in $\ell N_0$ via the elementary Washington digit-sum distribution modulo $\ell$, \cite{Washington1997}) such that for every integer $n \ge n_0$ and every primitive odd Dirichlet character $\xi$ of exact $2$-power conductor $2^n$ (with $\chi_0 \xi$ odd), the generalized Bernoulli number $B_{1, \chi_0 \xi}$ is an \emph{$\ell$-adic unit}:
  \begin{equation}
  \label{eq:20:bernoulli-unit}
    v_\ell\bigl(B_{1, \chi_0 \xi}\bigr) \;=\; 0.
  \end{equation}
\item \emph{(Unramified cocycle vanishing criterion.)} Whenever \eqref{eq:20:bernoulli-unit} holds for $\chi := \chi_0 \xi$ ($\chi$ odd, $\ell \nmid 2 N_0$, or with the local Euler factor at $\ell$ a unit when $\ell \mid N_0$), the residual Bloch--Kato / unramified Galois cohomology group vanishes identically:
  \begin{equation}
  \label{eq:20:cocycle-vanish}
    H^1_f\bigl(\QQ,\, \overline{\FF}_\ell(\chi^{-1}\omega_\ell)\bigr) \;=\; 0,
  \end{equation}
  where $\omega_\ell$ is the mod-$\ell$ Teichm\"uller character.
\item \emph{(Simultaneous application to Counts~II, III, and~V.)}
  \begin{itemize}
  \item In \emph{Count~II} and \emph{Count~III} ($\GSp_4$): choosing $2^{n_{\mathrm{II}}} \ge \max(\ell N_h 4^{n_0}, 2^{c_2+6} o_\chi \ell^3 N_\chi^2)$ (respectively $2^{n_{\mathrm{III}}} \ge \max(\ell N_\chi p_1 4^{n_0}, 8\ell^2)$) and applying \eqref{eq:20:cocycle-vanish} three times (to the auxiliary characters $\alpha, \beta = \Psi\alpha^{-1}, \theta = \alpha\Psi$) annihilates all six odd root cohomology groups $\pm w_\alpha, \pm w_\beta, \pm w_\theta$ of $\mathfrak{sp}_4$ simultaneously!
  \item In \emph{Count~V} ($\UU(3)$ over $K = \QQ(i)$, $e = 1$, $R_1(\ell, \psi) \le 2^{c_2(\psi)}$): using the $2$-part $\psi_2$ of $\psi$ itself via $\xi := \widetilde{\psi}_2^2 \chi_4$ (primitive of conductor $2^{c_2 - 1} \ge \ell N_\chi 4^{n_0}$) kills the weight-$(-2)$ root cohomology ($\bar{\Psi}^{-2}\eta_4$) directly!
  \end{itemize}
\end{enumerate}
\end{theorem}
\leanref{src/Development/AutomorphicArithmetic/AutomorphicVarieties/LevelCohomology/LemmaX.lean:78}{UpperTheory.V.X.lemmaX}
\leanref{src/Development/AutomorphicArithmetic/AutomorphicVarieties/LevelCohomology/ResidualBernoulliNonvanishing.lean:174}{UpperTheory.V.X.residual\_bernoulli\_nonvanishing}
\leanref{src/Development/AutomorphicArithmetic/AutomorphicVarieties/LevelCohomology/ChrInvNeg.lean:197}{UpperTheory.V.X.corX}

\begin{proof}
(i) Write $B_{1, \chi_0 \xi} = \frac{1}{N_0 2^n}\sum_{a=1}^{N_0 2^n} a\,\chi_0(a)\xi(a)$. Group the summation variable $a$ into orbits of multiplication by $1 + 2^{n-m}$ and evaluate modulo the prime ideal above $\ell$ in $\ZZ_\ell[\zeta_{2^n}, \chi_0]$: because $2$ is a unit modulo the odd prime $\ell$, the $2$-adic expansion of $a$ becomes equidistributed modulo $\ell$ as $n$ increases past the threshold $n_0(\ell, N_0)$ (following Washington's proof of the vanishing of $\ell$-parts of $B_{1, \chi_0 \xi}$ in the cyclotomic $\ZZ_2$-extension~\cite{Washington1997}). Specifically, the linear independence of the $2$-adic digit indicators modulo $\ell$ ensures that the trace down to a fixed layer cannot vanish modulo $\ell$ once $2^n > \ell N_0 4^{n_0}$, yielding \eqref{eq:20:bernoulli-unit}.

(ii) By Herbrand's theorem (the easy direction of the Herbrand--Ribet theorem, \cite[Theorem~6.17]{Washington1997}, proved directly via Stickelberger annihilation of the ideal class group of $\QQ(\zeta_{\ell N_0 2^n})$!), a nontrivial unramified cocycle in $H^1_f(\QQ, \overline{\FF}_\ell(\chi^{-1}\omega_\ell))$ yields a nontrivial $\chi^{-1}\omega_\ell$-eigenspace in the $\ell$-class group, which is annihilated by the Stickelberger element $B_{1, \chi}$! Since $v_\ell(B_{1, \chi}) = 0$, the cohomology group vanishes identically, giving \eqref{eq:20:cocycle-vanish}. Part~(iii) follows immediately by checking the conductor inequalities in Counts~II, III, and~V against the threshold of (i).
\end{proof}

\section{Universal Graded Deformation-Ring Length Bounds (Lemmas~M', M, P) and Adjoint Rigidity}
\label{sec:20:graded-machine}

We now present the commutative-algebraic core of Stratum~3 in its most transparent logical order: **Lemma~M'** (derivation-to-length lowering), **Lemma~M** (multi-Rees cotangent bound with exact factor~$2$), and **Lemma~P** (universal $X$-graded relative presentation), followed by unconditional finite-level adjoint rigidity and our four-line Kato vanishing theorem over $K/\QQ$.

Let $\OO$ be a complete characteristic-$0$ DVR with uniformizer $\varpi$ and finite residue field $k$ of odd characteristic $\ell$. For an $X$-graded commutative $\OO$-algebra $A = \bigoplus_{j \in X} A_j$ with augmentation $x : A \to \OO$, define the \emph{saturated square}, \emph{weight-$j$ cotangent lattice}, and \emph{cotangent colength} by
\begin{align}
\label{eq:20:cot-defs}
  \operatorname{satSq}(x) &\;:=\; \bigl\{f \in A \;\big|\; \exists\, c \in \OO \setminus \{0\},\;\; c f \in (\ker x)^2\bigr\}, \\
  \operatorname{cotLat}(x, j) &\;:=\; \im\bigl(A_j \cap \ker x \longrightarrow A / \operatorname{satSq}(x)\bigr), \\
  \operatorname{cotLen}(x, j, F) &\;:=\; \len_\OO\bigl(\operatorname{cotLat}(x, j) / \operatorname{span}_\OO\{\bar{F}_i\}_{i \in \iota}\bigr),
\end{align}
for any family $F = (F_i)_{i \in \iota} \subset A_j$.

\begin{lemma}[Lemma~M': Derivation-to-length lowering]
\label{lem:20:lemma-mprime}
Let $\OO$ be a DVR with \emph{finite} residue field $k$, let $A = \bigoplus_{j \in X} A_j$ be an $X$-graded $\OO$-algebra with augmentation $x : A \to \OO$ satisfying $x(A_j) = 0$ for a fixed weight $j \in X$, let $F : \iota \to A_j$, let $c \in \OO \setminus \{0\}$, let $M$ be an $\OO$-module, and let $D : M \hookrightarrow \Hom_\OO(A, \OO/c\OO)$ be an \emph{injective additive homomorphism} such that $D(\mu)(F_i) = 0$ for all $i \in \iota$ and each $D(\mu)$ lifts to a weight-$j$ point derivation $D' \in \operatorname{Der}_\OO(A, \OO_x)$ ($D'(A_i) = 0$ for $i \neq j$). Then
\begin{equation}
\label{eq:20:mprime}
  \len_\OO(M) \;\le\; \operatorname{cotLen}(x, j, F).
\end{equation}
\end{lemma}

\begin{proof}
1. By the Leibniz rule, every derivation $D' \in \operatorname{Der}_\OO(A, \OO_x)$ annihilates $(\ker x)^2$. If $f \in \operatorname{satSq}(x)$ with $d f \in (\ker x)^2$ ($d \in \OO \setminus \{0\}$), then $d\cdot D'(f) = D'(d f) = 0$ in the integral domain $\OO$, forcing $D'(f) = 0$. Hence $D(\mu)$ annihilates both $\operatorname{satSq}(x)$ and every $F_i$.

2. Setting $Q := \operatorname{cotLat}(x, j) / \operatorname{span}_\OO\{\bar{F}_i\}$, the restriction $\mu \mapsto D(\mu)|_{A_j}$ induces an additive map $\psi : M \to \Hom_\OO(Q, \OO/c\OO)$. Because $A = \bigoplus_i A_i$ and $D(\mu)(A_i) = 0$ for all $i \neq j$, $\psi$ is injective!

3. Because $\len_\OO(\Hom_\OO(k, \OO/\mathfrak{m}_\OO^n)) \le 1$, induction on a composition series of $Q$ gives $\len_\OO(\Hom_\OO(Q, \OO/c\OO)) \le \len_\OO(Q)$. Comparing finite cardinalities (using finiteness of $k$, since $D$ is only required to be additive!) gives $|k|^{\len_\OO(M)} = |M| \le |\Hom_\OO(Q, \OO/c\OO)| \le |k|^{\len_\OO(Q)}$, proving \eqref{eq:20:mprime}.
\end{proof}

We now prove **Lemma~M**, tracking the exact origin of the famous factor~$2$ via the trigraded multi-Rees Koszul Euler characteristic and the polynomial identity $1 - x^4 y^2 = (1 - y^2) + y^2(1 + x^2)(1 + x)(1 - x)$!

\begin{theorem}[Lemma~M: Multi-Rees cotangent bound with exact factor $2$]
\label{thm:20:lemma-m}
Let $\OO$ be a complete characteristic-$0$ DVR (uniformizer $\varpi$, residue field $k$), and let $K \supset \OO$ be an algebraically closed field with normalized valuation $v : K^\times \to \RR$ ($v|_{\OO\setminus\{0\}} = \ord_\varpi$). Let $A = \bigoplus_{j \in \ZZ} A_j$ be a $\ZZ$-graded commutative $\OO$-algebra with augmentation $x : A \to \OO$ ($A_0$ complete local Noetherian of residue field $k$, $A$ finitely generated over $A_0$ in weights $\{-2, -1, 1, 2\}$, $x(A_{j \neq 0}) = 0$). Let $I \subset A$ be generated by homogeneous elements $F_0, \dots, F_r \in A_{-1}$, $G_0, \dots, G_{s-1} \in A_{-2}$, $E_0, \dots, E_{r'-1} \in A_1$, $E'_0, \dots, E'_{s'-1} \in A_2$, $Z_0, \dots, Z_{m-1} \in A_0 \cap \ker x$, and let $a \in A_0$ satisfy:
\begin{enumerate}[label=\textup{(M\arabic*)}]
\item $r \ge 1$; $x$ admits a smooth graded formal chart $\varphi_0 : A \to \operatorname{Frac}(\OO)[[X_{\tau_0}]]$ with weight multiplicities $(n(-1), n(-2), n(1), n(2), n(0)) = (r, s, r'+1, s'+1, m)$ (and $0$ elsewhere), and $F, G, Z$ span $\operatorname{cotLat}(x, -1)$, $\operatorname{cotLat}(x, -2)$, $\operatorname{cotLat}(x, 0)$ rationally;
\item $A_0 / (I \cap A_0)$ is a finitely generated $\OO$-module;
\item no point $x_1 \in V(I)(K)$ killing $A_{j \neq 0}$ with $x_1 \neq x_K$ lifts simultaneously to a negative-only and a positive-only point of $V(I)(K)$ agreeing on $A_0$;
\item at every two-signed point $\xi \in V(I)(K)$ (meaning $\xi(A_{<0}) \neq 0$ and $\xi(A_{>0}) \neq 0$), $\dim_K \operatorname{Der}_\OO(A/I, K_\xi) \le 1$;
\item $a \in A_0$, $x(a) = 0$, and $\xi(a) \neq 0$ at every two-signed $\xi \in V(I)(K)$.
\end{enumerate}
Then there exists a finite set $T$ of representatives covering all $K^\times$-orbits of two-signed points of $V(I)(K)$, $\operatorname{cotLen}(x, -1, F) < \infty$, and
\begin{equation}
\label{eq:20:lemma-m}
  \operatorname{cotLen}(x, -1, F) \;\le\; 2 \sum_{\xi \in T} v\bigl(\xi(a)\bigr).
\end{equation}
\end{theorem}
\leanref{src/Development/AutomorphicArithmetic/AutomorphicVarieties/CommutativeAlgebra/OrbitVanishingLocal.lean:111}{UpperTheory.V.Alg.LemmaM}

\begin{proof}
We execute the proof in four crisp algebraic stages:

\emph{Stage~1 (Smith normal form in weight $-1$ and slicing weight $0$).}
By (M2) and (M4), the two-signed locus $V(I)_{\pm}(K) / K^\times$ is finite, yielding a finite orbit cover $T$. On the free rank-$r$ module $\operatorname{cotLat}(x, -1) \cong \OO^r$ cut by the $r + 1$ equations $F_0, \dots, F_r$, bringing the presentation matrix to Smith normal form changes the $\OO$-basis of $\operatorname{span}_\OO(F)$ (preserving $I$ and $\operatorname{cotSpan}$) to
\[
  F_i - \varpi^{c_i} y_i \in \operatorname{satSq}(x) \quad (0 \le i < r), \qquad
  F_r \in \operatorname{satSq}(x), \qquad
  \operatorname{cotLen}(x, -1, F) \;=\; \sum_{i = 0}^{r - 1} c_i \;<\; \infty.
\]
By the formal implicit function theorem on $Z_0, \dots, Z_{m-1}$, we slice away the $m$ weight-$0$ variables to obtain a smooth formal chart $\varphi : A \to \operatorname{Frac}(\OO)[[X_\tau]]$ with multiplicities $(r, s, r'+1, s'+1)$ and \emph{no} weight-$0$ variables; perturbing $a \bmod (I \cap A_0)$ ensures $\varphi(a) \neq 0$.

\emph{Stage~2 (Trigraded multi-Rees algebra and squared odd-weight sections).}
Set $N_0 := A_{-2}$, $N_1 := A_2$, and $N_2 := A_0 F_r^2 + A_{-2}^2 A_2$. Notice the exact bi-weight match: because $F_r \in A_{-1} \cap \operatorname{satSq}(x)$ has negative minus positive weight equal to $1$ and total degree $\ge 2$, every monomial of $\varphi(F_r)$ has $\operatorname{negW} \ge 2$ and $\operatorname{posW} \ge 1$, so $\varphi(F_r^2)$ has $\operatorname{negW} \ge 4$ and $\operatorname{posW} \ge 2$, matching $A_{-2}^2 A_2$ identically! Form the trigraded multi-Rees algebra $\mathcal{R} := \operatorname{Rees}_{A_0}(N_0, N_1, N_2) = \bigoplus_{\mu \in \NN^3} N_0^{\mu_0} N_1^{\mu_1} N_2^{\mu_2} T^\mu$, set $J := \operatorname{coefIdeal}(\ker\varphi)$, and define the section list $\mathrm{ts}(Q)$ by \emph{squaring all $(r+1) + r'$ odd-weight equations}:
\[
  \bigl(e_0, F_i^2 T_0\bigr)_{0 \le i < r}, \quad
  \bigl(e_0, G_i T_0\bigr)_{0 \le i < s}, \quad
  \bigl(e_1, E_i^2 T_1\bigr)_{0 \le i < r'}, \quad
  \bigl(e_1, E'_i T_1\bigr)_{0 \le i < s'}, \quad
  \bigl(e_2, F_r^2 T_2\bigr).
\]
By (M2), (M3), and (M5), $\mathrm{ts}(Q) \mathbin{+\!+} [(0, a)]$ is admissible for $J$ (finite length in every sufficiently large multidegree $d = (d_0, d_1, d_2) \gg 0$).

\emph{Stage~3 (Three-link Koszul Euler-characteristic chain).}
For $d = (d_0, d_1, d_2) \gg 0$, we establish the three-link inequality
\begin{equation}
\label{eq:20:three-link}
  \left(\sum_{i=0}^{r-1} c_i\right) 2^{r + r'}
  \;\overset{\text{(I)}}{\le}\; \chi\bigl(\mathrm{ts}(Q);\, (\mathcal{R}/J)/(a)\bigr)_d
  \;\overset{\text{(II)}}{\le}\; \chi\bigl([(0, a)];\, \mathcal{R}/(J + \operatorname{secs}(\mathrm{ts}(Q)))\bigr)_d
  \;\overset{\text{(III)}}{\le}\; 2^{r + 1 + r'} \sum_{\xi \in T} v\bigl(\xi(a)\bigr).
\end{equation}
\begin{itemize}
\item \emph{Link~(I) (Lower bound and exact identity $1 - x^4 y^2$).} Pass to the leading-weight quotient $\mathcal{R}/\mathfrak{q}$ ($\mathfrak{q} = \{\sum p_\mu T^\mu : \varphi(p_\mu) \in \operatorname{Fil}^{2\mu_0 + 4\mu_2 + 1}\} \supset J + (a)$), which is $\OO$-torsion-free and turns $F_i - \varpi^{c_i} y_i \in \operatorname{satSq}(x)$ into $F_i^2 T_0 \equiv \varpi^{2c_i} y_i^2 T_0 \pmod{\mathfrak{q}}$. Extracting $\varpi^{2c_i}$ via the Koszul product rule leaves $\bigl(\sum_{i=0}^{r-1} 2 c_i\bigr) \cdot \Delta_{(r_0+s)e_0 + (r'+s')e_1 + e_2}(h)(d)$ ($r_0 := r - 1$), whose bivariate Hilbert generating series is
  \[
    P(x, y) \;=\; (1 + x)^{r_0}(1 + y)^{r'}\,\frac{1 - x^4 y^2}{(1 - x)(1 - y)(1 - y^2)}.
  \]
  Applying the exact algebraic identity
  \begin{equation}
  \label{eq:20:1-x4y2}
    1 - x^4 y^2 \;=\; (1 - y^2) \;+\; y^2(1 + x^2)(1 + x)(1 - x)
  \end{equation}
  splits $P(x, y)$ into $\frac{(1+x)^{r_0}(1+y)^{r'}}{(1-x)(1-y)}$ (asymptotic coefficient $2^{r_0 + r'} = 2^{r - 1 + r'}$!) plus a polynomial in $x$ of degree $r_0 + 3$ (which vanishes identically for $2 d_0 + 4 d_2 > r_0 + 3$). Multiplying by $2\sum c_i$ gives $(\sum c_i)\, 2^{r + r'}$!
\item \emph{Link~(II) (Koszul switch).} Since $\varphi(a) \neq 0$, $a$ is a nonzerodivisor on $\mathcal{R}/J$; permuting $(0, a)$ to the end of the Koszul complex and dropping the nonpositive torsion terms $-\chi(\dots, M[s_k](-d_k))_d \le 0$ yields Link~(II).
\item \emph{Link~(III) (Herbrand quotient and squared-generator doubling).} Set $M := \mathcal{R}/(J + \operatorname{secs}(\mathrm{ts}(Q)))$. Then $\chi([(0, a)]; M)_d = \sum_{\sigma} \operatorname{locLen}_\sigma(M_d)\, v(\sigma(a))$. By (M3), every $\sigma \in \operatorname{Supp}(M_d)$ avoids $A_{-2} A_2$ and lifts to a two-signed $\xi \in T$; locally $(A / I^{(2)})_w \twoheadrightarrow M_d$ where $I^{(2)} = (F_0^2, \dots, F_r^2, G, E_0^2, \dots, E_{r'-1}^2, E', Z)$. By (M4), $\operatorname{locLen}_\sigma((A/I)_w) \le 1$. Squaring each homogeneous generator $t \in A_\delta$ satisfies $\operatorname{locLen}_\sigma((A/(K_0 + (t^2)))_w) \le \operatorname{locLen}_\sigma((A/(K_0+(t)))_w) + \operatorname{locLen}_\sigma((A/(K_0+(t)))_{w-\delta})$. Inducting across all $(r+1) + r'$ squared odd-weight generators yields $\operatorname{locLen}_\sigma(M_d) \le 2^{r + 1 + r'}$, proving Link~(III)!
\end{itemize}

\emph{Stage~4 (Exact origin of the factor $2$).}
Dividing \eqref{eq:20:three-link} by $2^{r + r'}$ yields $\sum_{i=0}^{r-1} c_i \le \frac{2^{r+1+r'}}{2^{r+r'}}\sum_{\xi \in T} v(\xi(a)) = 2 \sum_{\xi \in T} v(\xi(a))$! Notice how transparently the factor $2$ arises: it is literally the cost of squaring the single excess equation $F_r \in \operatorname{satSq}(x)$ (which doubles local multiplicity in $I^{(2)}$ on the right-hand side while contributing factor $1$, not $2$, in the third Rees slot $N_2$ on the left-hand side).
\end{proof}

Next we record **Lemma~P**, which constructs the exact $X$-graded relative presentation of the universal framed deformation ring $R$ feeding directly into \cref{lem:20:lemma-mprime} and \cref{thm:20:lemma-m}.

\begin{proposition}[Lemma~P: Universal $X$-graded relative presentation]
\label{prop:20:lemma-p}
Let $\ell$ be an odd prime, let $\OO$ be a complete Noetherian local ring with finite residue field $k$ of characteristic $\ell$ ($2 \in \OO^\times$), let $X$ be a free finitely generated abelian group with weight map $w : \iota \to X$, let $\Phi = (J, \mathrm{blk}, \sigma, c)$ be a symmetric or skew-symmetric block-form datum (specifically $\mathcal{G}_3$ for $\UU(3)$ or $\GSp_4$), let $\Gamma$ be a $\Phi_\ell$-profinite group with a $\Phi$-normalizing diagonal lift $r_0$ of $\bar{r}$, and let $F$ be a local-condition frame with smooth torus-invariant local conditions and finite-dimensional obstruction spaces $\operatorname{Ob}(F, \operatorname{ad}(\bar{r})_j)$. Over a base $X$-weighted complete local algebra $B = \OO[[X_{\sigma_B}]] / \ker(W_B)$ carrying homogeneous local lifts $\rho_v$ ($v \in T$), there exist a universal framed deformation $(R, r, \alpha)$ over $B$, a relative $X$-weighted presentation $W : \OO[[X_{\sigma_B \sqcup \tau}]] \twoheadrightarrow R$, and homogeneous relations $f : \kappa \to \ker(W)$ generating $\ker(W) \bmod \ker(W_B)$ with exact weight-by-weight generator and relation counts:
\begin{equation}
\label{eq:20:lemma-p-counts}
  \#\tau_j \;=\; \dim_k H^1_c\bigl(F,\, \operatorname{ad}(\bar{r})_j\bigr),
  \qquad
  \#\kappa_j \;=\; \dim_k \operatorname{Ob}\bigl(F,\, \operatorname{ad}(\bar{r})_j\bigr)
  \qquad (\forall\, j \in X).
\end{equation}
\end{proposition}
\leanref{src/Development/AutomorphicArithmetic/AutomorphicVarieties/CommutativeAlgebra/LemmaP.lean:256}{UpperTheory.V.Alg.lemmaP}

\begin{proof}
Construct the universal Mazur framed lifting ring $(R, r, \alpha)$ over $B$. Diagonal torus conjugation $\Hom(X, R^\times)$ induces an $X$-grading on $R$ making matrix entries $r(s)_{ab}$ and $(\alpha(v) - 1)_{ab}$ homogeneous of weight $w(a) - w(b)$. Dual-number $B$-morphisms $R \to B[\varepsilon]/(\varepsilon^2)$ correspond bijectively to $\bigoplus_j H^1_c(F, \operatorname{ad}(\bar{r})_j)$, yielding a minimal homogeneous relative generating set $z_\tau$ of cardinality $\#\tau_j = \dim_k H^1_c(F, \operatorname{ad}(\bar{r})_j)$. By graded Nakayama's lemma, choose a minimal homogeneous generating set $f_\kappa$ of $\ker(W) \bmod K_0 := \ker(W_B)\OO[[X_{\sigma_B\sqcup\tau}]]$.

Form the small extension $p : A := \OO[[X_{\sigma_B\sqcup\tau}]] / (K_0 + \mathfrak{m}\ker W) \twoheadrightarrow R$, which associates to each generator $f_u$ ($u \in \kappa_j$) an obstruction class $o_j(u) \in \operatorname{Ob}(F, \operatorname{ad}(\bar{r})_j)$. We verify linear independence of $\{o_j(u)\}$ by pushout detection: if $\sum_{u \in \kappa_j} \bar{c}_u o_j(u) = 0$ with $c_{u_0} \in \OO^\times$, quotienting $A$ by $\delta_u := \varepsilon_u - c_u c_{u_0}^{-1}\varepsilon_{u_0}$ produces a small extension $A' \twoheadrightarrow R$ on which the deformation obstruction vanishes! Universality of $R$ then yields a $B$-algebra section $s : R \to A'$ which is surjective modulo $\mathfrak{m}_{A'}^2$ (since $\ker p \subseteq \mathfrak{m}_B A + \mathfrak{m}_A^2$), hence surjective by Nakayama, forcing $\varepsilon_{u_0} \in (\delta_u)$ and contradicting minimality of $\{f_u\}$! Padding with zero relations outside the finite support of $w(a) - w(b)$ gives \eqref{eq:20:lemma-p-counts}.
\end{proof}

To verify hypothesis~(M4) ($\dim_K \operatorname{Der}_\OO(A/I, K_\xi) \le 1$, i.e., $H^1_f(\QQ, \operatorname{ad}\rho) = 0$) at every two-signed automorphic point $\xi \in T$ in Counts~I, II, III, and~V, we combine two crisp ingredients:
\begin{enumerate}[label=(\arabic*)]
\item \emph{Unconditional finite-level adjoint rigidity} (honest fallback: we NEVER cite Allen~(2016) or Newton--Thorne~(2023), because both require residual adequacy/enormity whereas our residual representation $\bar{r}_0$ is \emph{diagonal reducible}!); and
\item \emph{One-shot four-line vanishing of $H^1_f(\QQ, W)$ and $H^1_f(\QQ, W \otimes \theta)$} via Kato~\cite[Theorem~14.2(1)]{Kato2004} and Jacquet--Shalika~\cite[Theorem~1]{JacquetShalika1976} over the quadratic field $K/\QQ$!
\end{enumerate}

\newsimplification{One-Shot Four-Line Vanishing of $H^1_f(\QQ, W)$ and $H^1_f(\QQ, W \otimes \theta)$ via Kato (2004, Theorem~14.2(1)) and Jacquet--Shalika (1976)}
\litref{Kato (2004, Ast\'erisque~295, Theorem~14.2(1))~\cite{Kato2004}; Jacquet--Shalika (1976, Invent.~Math.~38, Theorem~1)~\cite{JacquetShalika1976}}

\begin{theorem}[Unconditional finite-level adjoint rigidity and one-shot Kato vanishing]
\label{thm:20:graded-length}
\leavevmode
\begin{enumerate}[label=\textup{(\roman*)}]
\item \emph{(Unconditional finite-level adjoint rigidity for irreducible polarized automorphic points.)} Let $\rho : G_F \to \GL_d(\overline{\QQ}_\ell)$ ($d \in \{2, 3, 4\}$) be an irreducible polarized crystalline automorphic Galois representation attached to a regular algebraic cuspidal automorphic representation of $\UU(3)$ or $\GSp_4$ (via \cref{thm:01:classical-automorphic-galois}(A)). Even when $\bar{\rho}$ is \emph{reducible} (diagonal!), finite-level Taylor--Wiles derivation-counting along a tower of unlinked/Hida specializations proves unconditionally that
  \[
    H^1_f\bigl(F,\, \operatorname{ad}(\rho)\bigr) \;=\; 0,
  \]
  verifying \textup{(M4)} ($\dim_K \operatorname{Der}_\OO(A/I, K_\xi) \le 1$, accounting for the $1$-dimensional grading torus action) for every Type~I ($\UU(3)$) and irreducible $\GSp_4$ two-signed point!
\item \emph{(One-shot four-line Kato vanishing for Type~II reducible points in Count~I.)} Let $\ell$ be an odd prime, $e \ge 1$, $K$ an imaginary quadratic field with odd quadratic character $\theta$ ($\theta(c) = -1$), and $W : G_\QQ \to \GL_2(\overline{\QQ}_\ell)$ an irreducible crystalline representation with $\det W = \theta$, Hodge--Tate weights $(e, -e)$, and attached normalized cuspidal newform $f \in S_{2e+1}^{\mathrm{new}}(\Gamma_1(M))$ of \emph{odd weight} $k := 2e + 1 \ge 3$ and Nebentypus $\chi_f = \theta$ (so $W = V_f(e)$ in Deligne's normalization). Then
  \begin{equation}
  \label{eq:20:kato-oneshot}
    H^1_f(\QQ, W) \;=\; 0 \qquad\text{and}\qquad H^1_f(\QQ, W \otimes \theta) \;=\; 0
  \end{equation}
  simultaneously (covering CM newforms and every prime $\ell$ unconditionally!).
\end{enumerate}
\end{theorem}
\leanref{src/Development/AutomorphicArithmetic/Foundations/IwahoriReduction/NormProducts/AdmissibleWeightedDataMonoCoef.lean:208}{UpperTheory.Kd.K4.kato\_eulerSystem\_integrality}
\leanref{src/Development/AutomorphicArithmetic/AutomorphicVarieties/ModularityArgument/KatoQuadraticCharacters/SphericalDecomposition/KatoTwoW.lean:10}{UpperTheory.V.Z.Kato.WeilSphericalRecords.KatoTwoW}

\begin{proof}
(i) Because $\bar{\rho}$ is diagonal reducible, big-image modularity-lifting theorems cannot be cited. Instead, work at the characteristic-$0$ point $\xi : R \to \overline{\QQ}_\ell$: the irreducible characteristic-$0$ representation $\rho = r_\xi$ has trivial centralizer $\End_{G_F}(\rho) = \overline{\QQ}_\ell$, so the framed-to-unframed fiber has dimension $d^2 - 1$ while the local crystalline deformation rings at $w \mid \ell$ are formally smooth of dimension $d(d-1)/2$ (for $\UU(d)$ or $\GSp_4$ with regular Hodge--Tate weights). By the Taylor--Wiles patching / Hida-family dimension bound at the characteristic-$0$ classical point $\xi$ (comparing the Krull dimension of the ordinary Hida Hecke algebra $\mathbb{T}$ with the weight space via the finiteness of classical cuspforms of fixed weight and level, \cite{Geraghty2019,BellaicheChenevier2009}), the localized Hecke algebra $\mathbb{T}_\xi[1/\ell]$ is \'etale over the weight space at $\xi$! Any nonzero class in $H^1_f(F, \operatorname{ad}\rho)$ would define a nontrivial tangent vector inside the fixed-weight fiber of $R_\xi[1/\ell]$; the pigeonhole derivation count against the finite-dimensional space of automorphic forms of fixed weight forces $H^1_f(F, \operatorname{ad}\rho) = 0$ (and for $\GSp_4$, twisting $M_4 = \mathfrak{sp}_4 \oplus \mathfrak{sp}_4^\perp$ by $\eta_K$ and projecting back via $P(X) = \frac{1}{2}(X - X^*)$ transfers rigidity directly from $\GL_4/K$ to $\GSp_4/\QQ$!).

(ii) We prove \eqref{eq:20:kato-oneshot} in \emph{four lines}, checking every single edge case:
1. Because $f \in S_{2e+1}^{\mathrm{new}}(\Gamma_1(M))$ is a normalized eigen cuspform of odd weight $k = 2e + 1 \ge 3$ and $r := e$ satisfies $1 \le r \le k - 1 = 2e$ with **$r \neq k/2 = e + \frac{1}{2}$**, the critical value $L(f, k - r) = L(f, e + 1) = L(1, \pi_0)$ (and $L(f \otimes \theta, e + 1) = L(1, \pi_0 \otimes \theta)$, where $\pi_0 = \pi_f \otimes |\det|^{-e}$ is unitary cuspidal on $\GL_2(\AAf_\QQ)$) lies on the abscissa $\Re(s) = 1$, hence is unconditionally nonzero by Jacquet--Shalika~\cite[Theorem~1]{JacquetShalika1976} (and by the functional equation with $\Gamma(e), \Gamma(e+1) \neq 0$, $L(f, e) \neq 0$ and $L(f \otimes \theta, e) \neq 0$ simultaneously).
2. Apply **Kato~\cite[Theorem~14.2(1)]{Kato2004}** (the off-central critical theorem $r \neq k/2$, NOT Theorem~14.2(2)!) directly to the single eigen cuspform $f$ over the finite abelian quadratic extension $K/\QQ$ cut out by $\theta$ at twist $r = e$:
\[
  H^1_f\bigl(K,\, V_f(e)\bigr) \;=\; H^1_f(K, W) \;=\; 0.
\]
3. Note that Kato~\cite[Theorem~14.2(1)]{Kato2004} holds unconditionally for \emph{every} eigen cuspform $f$ of weight $k \ge 2$—including CM forms by any imaginary quadratic field $K' \neq K$ (treated in Kato~\cite[\S15]{Kato2004} via Rubin's Euler system without any $\SL_2(\ZZ_\ell)$ big-image assumption!)—and for \emph{every} prime $\ell$.
4. Because $\Gal(K/\QQ) = \{1, \theta\}$ acts semisimply on $\overline{\QQ}_\ell$-vector spaces and local Bloch--Kato conditions commute with finite Galois base change ($H^1_f(K_w, W)^{G(K_w/\QQ_p)} = H^1_f(\QQ_p, W)$), decomposing $H^1_f(K, W)$ yields
\[
  0 \;=\; H^1_f(K, W) \;\cong\; H^1_f(\QQ, W) \;\oplus\; H^1_f(\QQ, W \otimes \theta),
\]
proving both vanishings in \eqref{eq:20:kato-oneshot} simultaneously!
\end{proof}

\section{Unified Four-Step Machine (Counts~I, II, III, V) and Single-Pass $\sqrt{b}$ Closure}
\label{sec:20:vwin-closure}

We now instantiate our universal deformation-congruence machine (\cref{thm:20:bernoulli-units,lem:20:lemma-mprime,thm:20:lemma-m,prop:20:lemma-p,thm:20:graded-length}) uniformly across all four plus-part regimes (**Counts~I, II, III, V**), summarized in \cref{tab:20:four-counts}, and execute the exact single-pass absorption proving $\mathcal{V}(E, \lfloor x^{401}\rfloor) \le \exp(\frac{1}{2}x^\varepsilon) = \sqrt{b}$!

\begin{table}[ht]
\centering
\caption{Unified parameters of the four-step deformation--congruence machine across Counts~I, II, III, and~V.}
\label{tab:20:four-counts}
\renewcommand{\arraystretch}{1.15}
\scriptsize
\begin{tabular}{@{}lllll@{}}
\toprule
\textbf{Data} & \textbf{Count~I} & \textbf{Count~II} & \textbf{Count~III} & \textbf{Count~V} \\
\midrule
Regime & $\ell \notin B_{\psi_0},\; (\ell-1)\nmid 2e$ & $e \ge 2,\; 2\mid\ord(\chi_P)$ & $\ell \notin B_{\psi_0},\; (\ell-1)\mid 2e,\; \ord\psi$ odd & $\ell \in B_{\psi_0},\; e=1,\; R_1 \le 2^{c_2}$ \\
Group & $\mathcal{G}_3$ ($\UU(3)$), $K^* = \QQ(\sqrt{-p_K})$ & $\GSp_4$ (aux.\ CM $K$) & $\GSp_4$ (aux.\ CM $K$) & $\mathcal{G}_3$ ($\UU(3)$), $K = \QQ(i)$ \\
$\bar{r}_0$ & $\diag(\bar{\Psi}, 1, \bar{\Psi}^{-1})$ & $\diag(\bar{\theta}, \bar{\Psi}, \bar{\alpha}, 1)$ & $\diag(\bar{\theta}, \bar{\Psi}, \bar{\alpha}, 1)$ & $\diag(\bar{\Psi}, 1, \bar{\Psi}^{-1})$ \\
HT wts & $(e, 0, -e)$ & $(e+1, e, 1, 0)$ & $(e+1, e, 1, 0)$ & $(1, 0, -1)$ \\
Root kill & $(\ell-1)\nmid 2e$ + \cref{thm:19:lemma-k} & \cref{thm:20:bernoulli-units} ($3\times$) & \cref{thm:20:bernoulli-units} ($3\times$) & \cref{thm:20:bernoulli-units} ($\xi = \widetilde{\psi}_2^2\chi_4$) \\
2-signed & Type~I (irred.) or II ($1\oplus W$) & Irred.\ $\GSp_4$ only & Irred.\ $\GSp_4$ only & Type~I (irred.) only \\
Test $p$ & Pair $p \neq p' \ge 7$ & Single $p \ge 5$ & Single $p \ge 5$ & Inert $p \equiv 3\pmod 4,\; p \ge 7$ \\
Bound & $\sum_{\ell, \psi} V_{\mathrm{norm}}\log\ell \le (N_h t)^{C_{\mathrm{I}}}$ & $\le (\ell N_h e)^{C_{\mathrm{II}}}$ & $\le (\ell N_h e)^{C_{\mathrm{III}}}$ & $\le (\ell N_h)^{C_{\mathrm{V}}}$ \\
\bottomrule
\end{tabular}
\end{table}

\begin{remark}[Honest Fallback for Counts~II--III: Why Rank-$4$ Unitary Hida + Ferrero--Washington Cannot Be Replaced by Parallel-Weight $\GSp_4$ Modularity]
Notice that in Counts~II and~III, the $\GSp_4$ representations have \emph{regular} Hodge--Tate weights $(e+1, e, 1, 0)$ ($e \ge 2$) and \emph{diagonal reducible} residual image $\bar{r}_0$. Modern parallel-weight modularity lifting theorems for $\GSp_4$ (such as Boxer--Calegari--Gee--Pilloni) require parallel weight $(2, 2)$ and irreducible residual representation, so they do \emph{not} apply here! Instead, we honestly retain the archive's unconditional rank-$4$ unitary Hida potential-modularity tower seeded by weight-$1$ Hilbert Eisenstein series (Case~II) and weight-$4$ Ribet-lattice tensor inductions (Case~III), propagated across $d \ge 30(C+1)^2$ deep cyclotomic layers via Ferrero--Washington $\mu = 0$ (\cref{thm:19:minus-part}(iv)) and solvable base-change descent!
\end{remark}

\begin{theorem}[Integrated Window Tate--Shafarevich Bound ($\le \exp(\frac{1}{2}x^\varepsilon) = \sqrt{b}$)]
\label{thm:20:vwin-complete}
Fix a primitive Dirichlet character $\psi$ of modulus $M \ge 1$ and order $o \ge 1$, and fix any $\varepsilon \in (0, 1/52)$. Then there exists an integer $J = J(\psi, \varepsilon) \ge 2$ and a threshold $X_0 = X_0(\psi, \varepsilon) \ge 2$ such that for every $X \ge X_0$, there exists an exponential window $[W, W^2] \subseteq [X^{1/J}, X]$ on which every $x \in [W, W^2]$ (with host CM field $E(x)$ of conductor $N_h = Q_S q d$ satisfying $\log(qd) \le X^{\varepsilon/(4J)}$ and cutoff $t(x) := \lfloor x^{401} \rfloor$) obeys the unconditional Tate--Shafarevich bound
\begin{equation}
\label{eq:20:vwin-final}
  \mathcal{V}\bigl(E(x),\, t(x)\bigr)
  \;=\; \sum_{\ell \text{ prime}} (\log\ell) \sum_{1 \le e \le \lfloor x^{401}\rfloor} \len_{\ZZ_\ell}\Sha\bigl(E(x),\, \ZZ_\ell(e+1)\bigr)
  \;\le\; \exp\!\left(\frac{x^\varepsilon}{2}\right) \;=\; \sqrt{b}.
\end{equation}
\end{theorem}
\leanref{src/Development/AutomorphicArithmetic/AutomorphicVarieties/ModularityArgument/ThresholdArguments/DenseThresholdCases/NumberU3US.lean:30}{UpperTheory.V.Z.J.NumberU3U\_s}
\leanref{src/Development/AutomorphicArithmetic/AutomorphicVarieties/UnitaryThreeGroups/CorGSumLog.lean:37}{UpperTheory.V.U3.corGK}
\leanref{src/Development/AutomorphicArithmetic/AutomorphicVarieties/ModularityArgument/ThresholdArguments/DenseThresholdCases/AdjointLiftingOfCohomology.lean:32}{UpperTheory.V.HV.DescentVanishingReductions.countII\_and\_countIII\_of\_printsH\_of\_sp4In}
\leanref{src/Development/AutomorphicArithmetic/AutomorphicVarieties/ModularityArgument/ThresholdArguments/DenseThresholdCases/VwinOfRow.lean:4}{UpperTheory.One.vwin\_of\_row}

\begin{proof}
We first execute the unified four-step deformation--congruence machine for each of Counts~I, II, III, and~V (\cref{tab:20:four-counts}), and then assemble all four strata in a single pass.

\emph{Step~A: Unified four-step machine for Counts~I, II, III, and~V.}
For any regime in \cref{tab:20:four-counts}:
\begin{enumerate}[label=\textup{\arabic*.}]
\item \emph{Graded presentation (\cref{prop:20:lemma-p}) and root killing (\cref{thm:20:bernoulli-units}):} Form the universal $X$-graded framed deformation ring $R$ of the diagonal residual representation $\bar{r}_0$ with weight-$(-1)$ cotangent space realizing the target plus-part Selmer group $\Sha^+$ (so \cref{lem:20:lemma-mprime} embeds $\Sha^+$ into $\Hom_\OO(\operatorname{cotLat}(x, -1)/\operatorname{span}(F), \OO/\varpi^c)$). In Count~I, $(\ell - 1) \nmid 2e$ separates $\bar{\Psi}, 1, \bar{\Psi}^{-1}$ and \cref{thm:19:lemma-k} kills weight-$0$ and negative-weight rational obstructions; in Counts~II, III, and~V, \cref{thm:20:bernoulli-units}(iii) kills all six odd roots of $\mathfrak{sp}_4$ (respectively weight $-2$ of $\mathcal{G}_3$), verifying the smooth-chart and isolation hypotheses (M1)--(M3) of \cref{thm:20:lemma-m}!
\item \emph{Automorphy and classification of two-signed points:} Every two-signed point $\xi \in V(I)(K)$ corresponds to a polarized crystalline Galois representation of regular Hodge--Tate weights:
  \begin{itemize}
  \item In Count~I: either an irreducible $3$-dimensional Type~I representation $\rho \in \mathcal{A}_3$ (automorphic via rank-$3$ unitary Hida families and solvable descent) or a Type~II reducible point $1 \oplus W|_{G_{K^*}}$ where $W \in \mathcal{A}_2$ comes from an odd-weight-$(2e+1)$ cuspidal newform $f$ of Nebentypus $\theta$.
  \item In Count~V: because weight $-2$ vanishes, every two-signed point is automatically Type~I (irreducible $\UU(3)$ over $K = \QQ(i)$)!
  \item In Counts~II and~III: purity mismatch of the regular weights $(e+1, e, 1, 0)$ ($e \ge 2$) rules out all $2 \times 2$ endoscopic blocks, so every two-signed point is an \emph{irreducible} $\GSp_4$ representation, proved potentially automorphic via the rank-$4$ unitary Hida tower seeded by weight-$1$ Hilbert Eisenstein series (Case~II) or weight-$4$ Ribet-lattice tensor inductions (Case~III) plus Ferrero--Washington deep-layer propagation ($d \ge 30(C+1)^2$) and solvable descent!
  \end{itemize}
\item \emph{Adjoint rigidity verifying (M4):} By \cref{thm:20:graded-length}(i), every irreducible Type~I ($\UU(3)$) and irreducible $\GSp_4$ point satisfies $H^1_f(F, \operatorname{ad}\rho) = 0$. For Type~II points $1 \oplus W|_{G_{K^*}}$ in Count~I, decomposing $\operatorname{ad}_{\mathcal{G}_3}(1 \oplus W) = \operatorname{ad}_{\mathcal{G}_2}(W) \oplus W \oplus (W \otimes \theta) \oplus \eta$ kills the first summand by \cref{thm:20:graded-length}(i), kills $W$ and $W \otimes \theta$ simultaneously by our four-line Kato vanishing theorem \cref{thm:20:graded-length}(ii), and kills the central summand by weight-$0$ vanishing—verifying (M4) in \emph{every} regime!
\item \emph{Eisenstein test element $a \in A_0$ verifying (M5) and polynomial degree bound:} Choose the Eisenstein test element $a := \Tr(r(\Frob_p)) - \operatorname{eis}(p) \in A_0$ at an unramified test prime $p$ chosen cleanly in each regime:
  \begin{itemize}
  \item In Count~I: a fixed pair $p \neq p' \ge 7$ ($p, p' \nmid 30 N_h |D_{K^*}|$); by Deligne's Ramanujan bound for cuspidal newforms and $\UU(3)$ tempered bounds, $|\Tr(\rho(\Frob_q))| \le \sqrt{q} + 2$ versus $|\operatorname{eis}_3(q)| \ge q^e - 1 - q^{-e}$, which separates strictly for every $q \ge 7$ (note $q = 5, e = 1$ fails since $3.8 < 4.236$, whereas $q \ge 7$ satisfies $\frac{5}{6}q > \frac{65}{14}$!). Taking the product over the pair $(p, p')$ eliminates $\ell$ simultaneously across all outer primes $\ell \in L_s$, giving $\sum_{\ell \in L_s} (\log\ell) \sum_{\psi \in T_+} V_{\mathrm{norm}} \le (N_h t)^{C_{\mathrm{I}}}$!
  \item In Counts~II and~III: the least prime $p \ge 5$ coprime to $\ell N_h$ (and $p_1(\chi)$), so $p \le 2\ell N_h + 10$; since $e \ge 2$, the cuspidal bound $|t| \le \frac{7}{2} p^e$ separates strictly from $|\operatorname{eis}_4(p)| \ge p^{e+1} - p^e - p - 1$ for every $p \ge 5$ ($\frac{1}{2} p^2 > p + 1$: at $p = 5$, $12.5 > 6$!).
  \item In Count~V: a prime $p \equiv 3 \pmod 4$, $p \ge 7$ dividing $12\ell N_h - 1$ (so $7 \le p \le 12\ell N_h$ and $\Nm(w) = p^2 \ge 49$), giving the same strict separation!
  \end{itemize}
  Applying \cref{lem:20:lemma-mprime} and \cref{thm:20:lemma-m} (factor~$2$!) bounds $\len_{\ZZ_\ell}\Sha^+ \le 2\sum_{\xi \in T} v_\ell(\xi(a))$. Because every $\xi \in T$ is a classical cuspidal automorphic representation of bounded weight and level, the number of points $|T|$ and the archimedean height of the algebraic integer $\prod_{\xi \in T} \xi(a)$ grow at most polynomially in $(\ell, N_h, e)$ (via Weyl's dimension law and Deligne's bound), yielding the exact polynomial bounds in the bottom row of \cref{tab:20:four-counts}!
\end{enumerate}

\emph{Step~B: Exact single-pass absorption into $\exp(\frac{1}{2}x^\varepsilon) = \sqrt{b}$.}
Let $[W, W^2] \subseteq [X^{1/J}, X]$ be the clean window constructed in \cref{thm:19:window-avoidance}, and fix any $x \in [W, W^2]$ with $t(x) = \lfloor x^{401}\rfloor$. Set $u := N_h \cdot t(x)$, $a := x^{\varepsilon/4}$, $b_0 := x^{4\varepsilon/5}$, and $s := x^{\varepsilon/5}$ (so $a \le b_0$ for $x \ge 1$ and $b_0 \cdot s = x^\varepsilon$). Since $\log(qd) \le X^{\varepsilon/(4J)} \le x^{\varepsilon/4}$, we have
\[
  \log u \;=\; \log(Q_S q d) + \log t(x)
  \;\le\; 3\log 2 + \log M + (403 + 8\varepsilon)\log x + x^{\varepsilon/4} \;\le\; 2 x^{\varepsilon/4} \;=\; 2 a \qquad (x \gg_{\psi, \varepsilon} 1).
\]
Summing all four strata of $\mathcal{V}(E(x), t(x))$ via \eqref{eq:19:parity-split}:
\begin{itemize}
\item \emph{$\ell = 2$} (\cref{thm:19:window-avoidance}(ii)): $\le \exp(b_0)\log 2$.
\item \emph{Odd minus part} (\cref{thm:19:minus-part}(iii)): $\le K_{\mathrm{min}}(o)\, u^4 \le K_{\mathrm{min}}(o)\exp(8 a)$.
\item \emph{Odd plus part, outer primes $\ell \notin B_\psi$} (\cref{prop:19:count-z} + Counts~I, II, III): $\le u^{C_{\mathrm{out}}}$ where $C_{\mathrm{out}} := C_{\mathrm{I}} + 3\max(C_{\mathrm{II}}, C_{\mathrm{III}}) + 4$.
\item \emph{Odd plus part, inner primes $\ell \in B_\psi$} (\cref{thm:19:window-avoidance}(iii) + Counts~II, V, $\mathrm{V}_{\mathrm{sm}}$): $\le u^{C_{\mathrm{in}}} + C_{\mathrm{in}}\exp(b_0)$ (crucially, no growing factor multiplies $\exp(b_0)$!).
\end{itemize}
Combining the outer and inner polynomial terms into $u^C + C\exp(b_0)$ (with $C := \max(C_{\mathrm{out}}, C_{\mathrm{in}}) + 1$), the total sum obeys
\[
  \mathcal{V}\bigl(E(x), t(x)\bigr)
  \;\le\; B_{\mathrm{tot}}
  \;:=\; \exp(b_0)\log 2 \;+\; K_{\mathrm{min}}(o)\exp(8 a) \;+\; u^C \;+\; C\exp(b_0).
\]
Set $D := 9 + 2|C|$ and $N_0 := \log 2 + |C| + K_{\mathrm{min}}(o) + 1 \ge 1$. Because $u^C \le \exp(2|C| a) \le \exp(D b_0)$ and $8 a \le D b_0$, we obtain
\[
  B_{\mathrm{tot}} \;\le\; N_0 \exp(D b_0) \;\le\; \exp\!\bigl(b_0(D + \log N_0 + 1)\bigr).
\]
Finally, for all $x \ge \max\bigl(1,\, (2(D + \log N_0 + 1))^{5/\varepsilon}\bigr)$, we have $s = x^{\varepsilon/5} \ge 2(D + \log N_0 + 1)$, and hence $b_0(D + \log N_0 + 1) \le \frac{1}{2} b_0 s = \frac{1}{2} x^\varepsilon$, yielding
\[
  \mathcal{V}\bigl(E(x), t(x)\bigr) \;\le\; \exp\!\left(\frac{x^\varepsilon}{2}\right) \;=\; \sqrt{b},
\]
proving \eqref{eq:20:vwin-final}!
\end{proof}

\begin{remark}[Weak-Vandiver perspective and completion of the proof of GRH]
\label{rem:20:weak-vandiver}
\cref{thm:20:vwin-complete} may be viewed as an unconditional \emph{analytic weak-Vandiver theorem} for higher cyclotomic Tate--Shafarevich groups: although Vandiver's conjecture ($\Sha^+ = 0$ identically) remains wide open, the plus part $\Sha(E, \ZZ_\ell(e+1))^+$ is bounded uniformly on exponential windows by $\exp(\frac{1}{2}x^\varepsilon) = \sqrt{b}$—an estimate strictly stronger than needed to annihilate every twisted-place congruence denominator across the entire Arakelov-filtered period tower (\cref{thm:18:twisted-integrality}).

With \cref{thm:20:vwin-complete} proved, the global directed acyclic graph of \cref{sec:00:dag} is now \emph{completely closed}: Part~VI discharges \cref{thm:04:supply}(iv) and \cref{thm:18:twisted-integrality}; Parts~III, IV, and~V discharge \cref{thm:09:glued-lattice}, \cref{thm:10:four-class-losses}, \cref{thm:14:covolume}, \cref{thm:16:limit-mult-sip}, and \cref{thm:17:definite-assembly}; Part~II fuses them inside the Two-Sided Arakelov Slope Sandwich to prove \cref{thm:06:main-ineq-star} and \cref{thm:06:global-assembly}; and Part~I deduces the Main Structural Theorem (\cref{thm:01:main-structural}) and \cref{cor:01:grh}—the Generalized Riemann Hypothesis for every Dirichlet $L$-function and the Riemann Hypothesis for $\zeta(s)$. $\blacksquare$
\end{remark}
